%% 
%% Copyright 2007-2026 Elsevier Ltd
%% 
%% This file is part of the 'Elsarticle Bundle'.
%% ---------------------------------------------
%% 
%% It may be distributed under the conditions of the LaTeX Project Public
%% License, either version 1.3 of this license or (at your option) any
%% later version.  The latest version of this license is in
%%    http://www.latex-project.org/lppl.txt
%% and version 1.3 or later is part of all distributions of LaTeX
%% version 1999/12/01 or later.
%% 
%% The list of all files belonging to the 'Elsarticle Bundle' is
%% given in the file `manifest.txt'.
%% 
%% Template article for Elsevier's document class `elsarticle'
%% with harvard style bibliographic references

\documentclass[preprint,11pt,authoryear]{elsarticle}

%% Use the option review to obtain double line spacing
%% \documentclass[authoryear,preprint,review,12pt]{elsarticle}

%% Use the options 1p,twocolumn; 3p; 3p,twocolumn; 5p; or 5p,twocolumn
%% for a journal layout:
%% \documentclass[final,1p,times,authoryear]{elsarticle}
%% \documentclass[final,1p,times,twocolumn,authoryear]{elsarticle}
%% \documentclass[final,3p,times,authoryear]{elsarticle}
%% \documentclass[final,3p,times,twocolumn,authoryear]{elsarticle}
%% \documentclass[final,5p,times,authoryear]{elsarticle}
%% \documentclass[final,5p,times,twocolumn,authoryear]{elsarticle}

%% For including figures, graphicx.sty has been loaded in
%% elsarticle.cls. If you prefer to use the old commands
%% please give \usepackage{epsfig}


\usepackage{amsmath,amssymb,amsthm}
\usepackage{booktabs}
\usepackage{hyperref}
\usepackage[margin=1in]{geometry}
\usepackage[utf8]{inputenc}
\usepackage{pdflscape}
\usepackage{tabularx}
\usepackage{caption}
\usepackage{setspace}
\onehalfspacing

\journal{European Journal of Operations Research}

\begin{document}

\begin{frontmatter}

%% Title, authors and addresses

%% use the tnoteref command within \title for footnotes;
%% use the tnotetext command for theassociated footnote;
%% use the fnref command within \author or \affiliation for footnotes;
%% use the fntext command for theassociated footnote;
%% use the corref command within \author for corresponding author footnotes;
%% use the cortext command for theassociated footnote;
%% use the ead command for the email address,
%% and the form \ead[url] for the home page:
%% \title{Title\tnoteref{label1}}
%% \tnotetext[label1]{}
%% \author{Name\corref{cor1}\fnref{label2}}
%% \ead{email address}
%% \ead[url]{home page}
%% \fntext[label2]{}
%% \cortext[cor1]{}
%% \affiliation{organization={},
%%            addressline={}, 
%%            city={},
%%            postcode={}, 
%%            state={},
%%            country={}}
%% \fntext[label3]{}

\title{Title Placeholder} %% Article title

%% use optional labels to link authors explicitly to addresses:
%% \author[label1,label2]{}
%% \affiliation[label1]{organization={},
%%             addressline={},
%%             city={},
%%             postcode={},
%%             state={},
%%             country={}}
%%
%% \affiliation[label2]{organization={},
%%             addressline={},
%%             city={},
%%             postcode={},
%%             state={},
%%             country={}}

\author{Author Name 1} %% Author name

%% Author affiliation
\affiliation{organization={Chair of Computational Logistics, RWTH Aachen University},%Department and Organization
            addressline={Kackerstra\ss{}e 7}, 
            city={Aachen},
            postcode={52072}, 
            state={NRW},
            country={Germany}}

%% Abstract
\begin{abstract}
%% Text of abstract
Abstract text.
\end{abstract}

%%Graphical abstract
\begin{graphicalabstract}
%\includegraphics{grabs}
\end{graphicalabstract}

%%Research highlights
\begin{highlights}
\item Research highlight 1
\item Research highlight 2
\end{highlights}

%% Keywords
\begin{keyword}
%% keywords here, in the form: keyword \sep keyword

%% PACS codes here, in the form: \PACS code \sep code

%% MSC codes here, in the form: \MSC code \sep code
%% or \MSC[2008] code \sep code (2000 is the default)

\end{keyword}

\end{frontmatter}

%% Add \usepackage{lineno} before \begin{document} and uncomment 
%% following line to enable line numbers
%% \linenumbers

%% main text
%%

%% Use \section commands to start a section
\section{Literature Review}

\paragraph{\textbf{The Robust Pickup and Delivery Problem with Time Windows}}

\citet{abreuRobustPickupDelivery2026} study the robust pickup and delivery
problem with time windows (RPDPTW), which handles demand and travel time uncertainty via robust
optimization rather than being assumed away, and propose both a compact robust counterpart formulation and a branch-and-cut algorithm.

The main contribution of the paper is to introduce the robust PDPTW to the literature.
It shows how PDPTW can be extended to handled uncertainty.


\paragraph{\textbf{An Exact Algorithm for the Pickup and Delivery Problem with Time Windows and Last-in-First-out Loading}}

\citet{alyasiryExactAlgorithmPickup2019} study the PDPTW with last-in-first-out loading constraints. 
Their main contribution is to propose a new exact method to solve PDPTWs, including the PDPTWL. Their solution method is based on 
fragments, which are series of pickup and delivery requests starting and ending with an empty vehicle. They use such fragments to 
create a relaxed network flow model with side constraints.

This paper is relevant for two reasons: it studies a PDPTW variant with LIFO, that is, a specific loading constraint, adjacent to our 
loading constraint (First In, All Out), and the solution method is somewhat interesting because their "fragments" are a similar concept to our 
delivery batches (sequence of pickup and delivery events starting and ending with an empty vehicle).


\paragraph{\textbf{An Exact Algorithm for the Pickup and Delivery Problem with Time Windows}}

\citet{baldacciExactAlgorithmPickup2011} study the standard PDPTW and develop a new exact algorithm based on set-partitioning formulation. 
They showcase a bounding procedure that finds near-optimal dual solutions to the linear relaxation, incorporating two dual ascent heuristics and cut-and-column 
generation procedure. The final dual solution is used to generate a reduced master problem, which is solved either by an integer programming solver or a branch-and-price algorithm, depending on size.

Their main contribution is the development of a new exact PDPTW algorithm that outperforms existing exact algorithms.

\paragraph{\textbf{The Pickup and Delivery Problem with Time Windows and Scheduling on the Edges}}

\citet{barbosaPickupDeliveryProblem2026} study the pickup and delivery problem with time windows
and scheduling on the edges (PDPTW-SE), which extends the PDPTW by allowing pickup and delivery
nodes to be located in different regions inaccessible to regular vehicles. 
For such requests, the vehicle must be transported across a special edge by a separate machine (e.g., ferry, elevator), and since only a limited number of machines are available, 
their use must be scheduled subject to time windows.

Their main contribution is to introduce the PDPTW-SE to the literature, which has applications
in settings such as the logistical planning of pickups and deliveries on multiple islands or AGV-based pickup-and-delivery planning in multi-floor buildings such as hotels or hospitals.


\paragraph{\textbf{Static Pickup and Delivery Problems: A Classification Scheme and Survey}}

\citet{berbegliaStaticPickupDelivery2007} introduce a general modeling framework and three-field classification scheme
covering the broad family of static PDPs, and survey the solution methods applied to each class.

Their taxonomy (one-to-one, one-to-many-to-one, many-to-many) is standard and can be used to situate our specific model relative to existing problem classes.


\paragraph{\textbf{A branch-and-price algorithm for the multi-depot heterogeneous-fleet pickup and delivery problem with soft time windows}}

\citet{bettinelliBranchandpriceAlgorithmMultidepot2014} study a PDPTW variant that addresses multiple depots, heteregenous fleets, and flexible preference handling simultaneously. This variant arises from 
growing needs and demands logistics managements need to fulfill. Their solution method is a branch-and-price algorithm.

This paper is relevant for two reasons: first, it addresses a multi-depot PDPTW variant, which is a deviation from the classic single-depot feature of the PDPTW, and our model does something similar (remove depots all together). 
Second, it uses a branch-and-price algorithm to solve this problem, which matches our solution method.


\paragraph{\textbf{Branch-Price-and-Cut Algorithms for the Pickup and Delivery Problem with Time Windows and Last-in-First-Out Loading}}

\citet{cherkeslyBranchPriceandCutAlgorithmsPickup2015} study the PDPTWL (PDPTW with Last-In-First-Out loading constraints), the first paper to develop an exact method to this problem variant. 
They propose three BCP algorithms, one in which the LIFO constraints are incorporated into the master problem, one in which they are included in the pricing problem, and a hybrid.

Mainly the introduction of LIFO constraints to the PDPTW is an interesting contribution relevant to our model.


\paragraph{\textbf{A Branch-and-Cut Algorithm for the Dial-a-Ride Problem}}

\citet{cordeauBranchandCutAlgorithmDialaRide2006} study the DARP, in which vehicles must transport passengers between pickup and delivery locations,
subject to time windows and maximum ride-time constraints. The DARP is a fairly widely studied PDPTW variant.

Their contribution is a branch-and-cut algorithm for the DARP, which incorporates known and new valid inequalities.


\paragraph{\textbf{A New Regret Insertion Heuristic for Solving Large-Scale Dial-a-Ride Problems with Time Windows}}

\citet{dianaNewRegretInsertion2004} study a large-scale
dial-a-ride problem with time windows (DARPTW).

Their main contribution is a parallel regret insertion heuristic for the
DARP. A regret-based insertion heuristic is a construction heuristic that
iteratively inserts requests into a partial solution in order of their regret
value, which is the difference in cost between the best and second-best
feasible insertion positions for that request. Though it is not super relevant for a branch-and-price framework, Sjaak's solver used the methodology to create a regret construction heuristic for his solvers.


\paragraph{\textbf{The Pickup and Delivery Problem with Time Windows}}

\citet{dumasPickupDeliveryProblem1991} study the standard pickup and
delivery problem with time windows (PDPTW).

The paper's main contribution is the first exact solution algorithm for
the PDPTW: a set-partitioning formulation whose LP relaxation is solved
by column generation, with columns priced by a constrained shortest-path
subproblem and a branch-and-bound tree used to resolve fractional
solutions by branching on pickup sequences.


\paragraph{\textbf{Pickup and Delivery Problem with Time Windows: A New Compact Two-Index Formulation}}

\citet{furtadoPickupDeliveryProblem2017} study the standard PDPTW and propose a new compact two-index formulation.

Their main contribution is the implementation of a new modeling strategy that allows the assignment of vehicles
to routes explicitly in two-index flow formulations. It has a higher performance relative to classical compact three-index formulations.


\paragraph{\textbf{Branch-and-Price for the Pickup and Delivery Problem with Time Windows and Scheduled Lines}}

\citet{ghilasBranchandPricePickupDelivery2018} propose a branch-and-price algorithm for the PDPTW-SL, a variant they previously introduced. The PDPTW-SL integrates freight flows 
with scheduled public transportation services. Vehicles are allowed to use scheduled public transportation lines to perform pickups and deliveries, which can reduce costs and improve service levels.

The main relevant contribution of this paper is a semi-recent application of a branch-and-price framework for a PDPTW variant. Additionally, 
the PDPTW-SL is somewhat relevant to us as another example of a restricted PDPTW variant in which the freedom to route between arbitrary origins and 
destinations is restricted by a structural constraint (fixed public-transport schedules). 


\paragraph{\textbf{Bidirectional Labeling in Column-Generation Algorithms for Pickup-and-Delivery Problems}}

\citet{gschwindBidirectionalLabelingColumngeneration2018} study the PDPTW, in particular within the context of branch-cut-and-price frameworks, 
and build upon the work of \citet{ropkeBranchCutPrice2009} to develop a bidirectional labeling algorithm for the pricing subproblem.
Specifically, strong dominance is applied in both forward and backward directions simultaneously, which allows for a more efficient exploration of the search space.
They demonstrate that this bidirectional strategy reduces the pricing time compared to forward labeling.

Their main contribution is the theoretical development of a bidirectional labeling algorithm and the empirical demonstration of 
its effectiveness in pricing time reduction for the PDPTW. Potentially, our model could implement bidirectional labeling, making this paper very relevant.


\paragraph{\textbf{A Fast Decomposition and Reconstruction Framework for the Pickup and Delivery Problem With Time Windows and LIFO Loading}}

\citet{liuFastDecompositionReconstruction2019} study the PDPTWL, a PDPTW variant incorporating LIFO loading. Their main contribution is to apply a Decomposition and Reconstruction algorithm coupled with a
tabu search to the PDPTWL, finding new BKS to many Li \& Lim instances for this specific variant.


\paragraph{\textbf{The pickup and delivery problem with transshipments: Critical review of two existing models and a new formulation}}

\citet{lyuPickupDeliveryProblem2023} study the PDP-T and PDPTW-T, which are extensions from the standard PDPTW to include transshipments, 
where requests can be transfered between vehicles at desginated transfer stations.

The paper's main contribution is to provide review and revision of the two existing state-of-the-art models for the PDP-T and PDPTW-T, 
while also developing a new mixed-integer programming formulation for these variants, which has better computational performance.


\paragraph{\textbf{An Adaptive Large Neighborhood Search Heuristic for the Pickup and Delivery Problem with Time Windows}}

\citet{ropkeAdaptiveLargeNeighborhood2006} study the PDPTW and extend Shaw's large neighborhood search (LNS) heuristic. 
Their main contribution is the development of an adaptive LNS heuristic that dynamically adjusts the selection of removal and insertion operators based on their past performance, 
which allows for a more effective exploration of the solution space and improved solution quality. Their repair operators include basic greedy heuristics as well as regret heuristics. 


\paragraph{\textbf{Branch and Cut and Price for the Pickup and Delivery Problem with Time Windows}}

\citet{ropkeBranchCutPrice2009} study the PDPTW and develop a BCP algorithm. Their main contribution is the introduction of valid inequalities in the formulation 
to improve lower bound quality. Another contribution is that they showcased how previously proposed inequalities for the PDPTW are implied by the set partitioning formulation.


\paragraph{\textbf{Models and Branch-and-Cut Algorithms for Pickup and Delivery Problems with Time Windows}}

\citet{ropkeModelsBranchandcutAlgorithms2007} study both the standard PDPTW and the DARP. They propose two new formulations 
for the PDPTW using branch-and-cut algorithms, and also introduce new valid inequalities.


\paragraph{\textbf{A Study on the Pickup and Delivery Problem with Time Windows: Matheuristics and New Instances}}

\citet{sartoriStudyPickupDelivery2020} study the standard PDPTW and propose a new matheuristic approach. Their main contribution, aside from the integration of their mathematical programming component, is the creation of extensive 
new routing benchmark instances based on real-world data. In our research, we use a modified version of their benchmark instances.


\paragraph{\textbf{The General Pickup and Delivery Problem}}

\citet{savelsberghGeneralPickupDelivery1995} study the
general pickup and delivery problem (PDP). The difference to the standard PDPTW
is that the PDP does not include time windows (and unlike 1D-PDPTW, no
single-destination requirement).

Their main contribution is definitional and survey-oriented rather than
algorithmic: they formalize the PDPTW problem definition and catalogue the
variety of objective functions used in the literature (e.g., minimizing fleet
size, total travel cost, total duration, or combinations thereof), which we
draw on directly when positioning our own cost-minimization objective (or some weighted variant).


\section{State Transitions Overview}

\begin{landscape}
\vspace*{\fill}
\noindent
\begin{minipage}{\linewidth}
\centering
\small
\renewcommand{\arraystretch}{1.5}
\setlength{\tabcolsep}{5pt}
\begin{tabularx}{\linewidth}{|>{\raggedright\arraybackslash}p{3.1cm}|
                                >{\raggedright\arraybackslash}X|
                                >{\raggedright\arraybackslash}X|
                                >{\raggedright\arraybackslash}X|
                                >{\raggedright\arraybackslash}X|
                                >{\raggedright\arraybackslash}X|}
\hline
\textbf{State attribute} &
\textbf{T1: Route start (first pickup)} &
\textbf{T2: Pickup $\rightarrow$ Route closure} &
\textbf{T3: Pickup $\rightarrow$ pickup, different destination} &
\textbf{T4: Pickup $\rightarrow$ pickup, same destination, no intermediate delivery needed} &
\textbf{T5: Pickup $\rightarrow$ pickup, same destination, intermediate delivery needed} \\
\hline\hline
Current node $v$ &
The route begins at the first pickup node $p(k)$, stored as $\pi_1$. &
The route terminates after final delivery. &
Route continues after delivery to the new pickup $p(k)$. &
The next node is the new pickup $p(k)$. &
Route continues after delivery to the new pickup $p(k)$. \\
\hline
Load $L$ &
$L' = q_k$. &
$L' = 0$. &
$L' = q_k$. &
$L' = L + q_k$. &
$L' = q_k$. \\
\hline
Time $T$ &
$T' = a_k^P + s_k^P$. &
$T' = T_\delta$. &
$T' = \max\{T_\delta + t_{\delta,p(k)},\,a_k^P\} + s_k^P$. &
$T' = \max\{T + t_{v,p(k)}, a_k^P\} + s_k^P$. &
$T' = \max\{T_\delta + t_{\delta,p(k)},\,a_k^P\} + s_k^P$. \\
\hline
Current cost $C$ &
$C' = 0$. &
$C' = C_\delta + \rho \cdot c_{\delta,\pi_1}$. &
$C' = C_\delta + c_{\delta,p(k)}$. &
$C' = C + c_{v,p(k)}$. &
$C' = C_\delta + c_{\delta,p(k)}$. \\
\hline
Current batch $o$ &
$o' = \{k\}$. &
$o' = \varnothing$. &
$o' = \{k\}$. &
$o' = o \cup \{k\}$. &
$o' = \{k\}$. \\
\hline
Initiated requests $S$ &
$S' = S \cup \{k\}$. &
$S' = S$. &
$S' = S \cup \{k\}$. &
$S' = S \cup \{k\}$. &
$S' = S \cup \{k\}$. \\
\hline
Active destination $\delta$ &
$\delta' = d(k)$. &
$\delta' = \varnothing$. &
$\delta' = d(k)$. &
$\delta' = \delta$. &
$\delta' = d(k)$. \\
\hline
Partial route $\tilde{\pi}$ &
Initialize with $p(k)$. &
Append $\delta$. &
Append $\delta$, $p(k)$. &
Append $p(k)$. &
Append $\delta$, $p(k)$. \\
\hline
\end{tabularx}
\captionof{table}{State transitions in the 1D-PDPTW state-space algorithm. The state attributes $v$, $L$, $T$, $C$, $o$, $S$, $\delta$, and $\tilde{\pi}$ are the attributes of a given state (where applicable) before the transition. The quantities $T_\delta$ and $C_\delta$ denote the output of $\textsc{Deliver}(o)$ at the active destination $\delta$. Deliveries are embedded in the transition logic where applicable. $\pi_1$ denotes the first physical node visited by the route, fixed at the start of the route (T1) and carried unchanged by all subsequent transitions; it is required to evaluate the non-loop penalty at route closure (T2). $\rho$ denotes the non-loop route penalty ratio. The state snapshot is taken after the vehicle completes service at the current node and is ready to depart.}
\label{tab:state-transitions}
\end{minipage}
\vspace*{\fill}
\end{landscape}


\section{Results}

\begin{table}[htbp]
\centering
\caption{Comparative Results of Stefan's branch-and-price solver against Sjaak's extended solver.}

\begin{tabular}{lrrrr}
    \toprule
    Instance & Stefan (obj) & Stefan time [s] & Sjaak (obj) & Sjaak time [s] \\
    \midrule

    bar-n100-1\_50\% & 663 & 0.62 & 666$^*$ & 69068.6 \\
    bar-n100-1\_75\% & 646.5 & 5.50 & 674$^*$ & 21998.3 \\
    bar-n100-2\_20req\_75\% & 297 & 0.01 & 297 & 3001.4 \\
    bar-n100-5\_25req\_50\% & 385 & 0.02 & -- & -- \\
    ber-n100-1\_20req\_25\% & 759.5 & 0.00 & -- & -- \\
    ber-n100-1\_25req\_25\% & 899 & 0.00 & -- & -- \\
    ber-n100-5\_50\% & 1402.5 & 0.05 & -- & -- \\
    ber-n100-5\_75\% & 1216 & 0.09 & -- & -- \\
    nyc-n100-3\_5\% & 111.5 & 47.9 & -- & -- \\
    nyc-n100-3\_25\% & 455 & 0.50 & -- & -- \\
    nyc-n100-3\_50\% & 536.5 & 0.36 & 540$^*$ & 22051.3 \\
    nyc-n100-3\_75\% & 519 & 0.33 & -- & -- \\
    nyc-n100-4\_25req\_5\% & 284 & 0.04 & -- & -- \\
    \bottomrule
\end{tabular}

\vspace{4pt}
\footnotesize
$^*$ Route enumeration for the extended model hit its configured time limit (3,000s / 21,000s / 68,400s) before completing all relevant request-subset sizes, so the reported objective and its 0\% MIP gap are only optimal with respect to the incomplete route pool at that point, not globally.

\end{table}
%% For citations use: 
%%       \citet{<label>} ==> Lamport (1994)
%%       \citep{<label>} ==> (Lamport, 1994)
%%

%% If you have bib database file and want bibtex to generate the
%% bibitems, please use
%%
%%  \bibliographystyle{elsarticle-harv} 
%%  \bibliography{<your bibdatabase>}

%% else use the following coding to input the bibitems directly in the
%% TeX file.

%% Refer following link for more details about bibliography and citations.
%% https://en.wikibooks.org/wiki/LaTeX/Bibliography_Management

\bibliographystyle{apalike-ejor}
\bibliography{references}
\end{document}

\endinput
%%
%% End of file `elsarticle-template-harv.tex'.


